\appendix

\chapter{Derivazione della distribuzione di Boltzmann}
\label{sec:a_boltzm_der}
Il comportamento del sistema a fissata $T$ è descritto da una distribuzione di probabilità $p (\mathbf{x}, T)$ ottenibile attraverso il principio di massima entropia.
\begin{equation}
    \mathcal{S} = - \int_{\mathbb{R}^D} p (\mathbf{x}, T) \ln(p(\mathbf{x}, T)) d\mathbf{x}
\end{equation}

Per massimizzare l'entropia si utilizza il teorema dei moltiplicatori di Lagrange:

\begin{equation}
    \mathcal{L}[p] = -\int_{\mathbb{R}^D} p \ln p \, d\mathbf{x} - \alpha \left( \int_{\mathbb{R}^D} p \, d\mathbf{x} - 1 \right) - \beta \left( \int_{\mathbb{R}^D} p \, \mathcal{H} \, d\mathbf{x} - \langle E \rangle \right),
\end{equation}

dove $\alpha$ e $\beta$ sono i moltiplicatori di Lagrange associati rispettivamente al vincolo di normalizzazione

\begin{equation}
    \int_{\mathbb{R}^D} p(\mathbf{x}, T) \ d\mathbf{x} = 1 \ ,
\end{equation}

e al vincolo sul valore medio dell'energia

\begin{equation}
    \int_{\mathbb{R}^D} p(\mathbf{x}, T) \, \mathcal{H}(\mathbf{x}) \, d\mathbf{x} = \langle E \rangle .
\end{equation}

La massimizzazione impone $\delta \mathcal{L}/\delta p = 0$ restituendo:

\begin{equation}
    - \ln p(\mathbf{x}, T) - 1 - \alpha - \beta \mathcal{H}(\mathbf{x}) = 0
\end{equation}

che permette di determinare $p (\mathbf{x}, T)$:

\begin{equation}
    p(\mathbf{x}, T) = e^{-1-\alpha} \, e^{-\beta \mathcal{H}(\mathbf{x})}.
\end{equation}

Il moltiplicatore $\alpha$ è fissato dalla condizione di normalizzazione: definendo la \textbf{funzione di partizione}

\begin{equation}
    Z(\beta) = \int_{\mathbb{R}^D} e^{-\beta \mathcal{H}(\mathbf{x})} \ d\mathbf{x},
\end{equation}

si ha $e^{-1-\alpha} = 1/Z(\beta)$, da cui \textbf{la distribuzione di Boltzmann}:

\begin{equation}
    p(\mathbf{x}, T) = \frac{1}{Z(\beta)} e^{-\beta \mathcal{H}(\mathbf{x})}.
\end{equation}

Dove $\beta = 1/(k_B T)$.
